% Irodov Part VI §6.1 problems 19--32.
% Original source printed pp.235--236, PDF pp.244--245.
% Distinguish principal quantum number n from ordinal line number:
% Balmer line k corresponds to upper n=k+2. This matters for #29.
\atomicproblem{19}{235}
用前题的经典辐射公式估计：一个电子沿半径
$r_0=50\,\mathrm{pm}$ 的氢原子圆轨道运动，
因连续辐射损失能量而落到核上需多久？
假定加速度始终指向核中心。
\exsolution
记库仑常数 $k_e=1/(4\pi\varepsilon_0)$。
电子圆轨道有
\[
 m_ev^2/r=k_ee^2/r^2,\qquad
 E(r)=\tfrac12m_ev^2-k_ee^2/r
 =-\frac{k_ee^2}{2r},
\]
并有 $a=k_ee^2/(m_er^2)$。
SI 制拉莫尔辐射功率
$P=e^2a^2/(6\pi\varepsilon_0c^3)
=2k_e^3e^6/(3m_e^2c^3r^4)$。
能量损失
$(dE/dr)(dr/dt)=-P$，
得到
\[
 \frac{dr}{dt}=-\frac{4k_e^2e^4}{3m_e^2c^3r^2}.
\]
从 $r_0$ 积分到0：
\[
 \boxed{t=\frac{m_e^2c^3r_0^3}{4k_e^2e^4}
 \simeq1.3\times10^{-11}\,\mathrm s=13\,\mathrm{ps}}.
\]
\excheck 这是古典模型的原子坍缩时间，
正是其无法描述真实稳定原子的经典矛盾。

\atomicproblem{20}{235}
证明：类氢离子电子从第 $n+1$ 轨道跃迁到第
$n$ 轨道时，辐射角频率 $\omega$
满足 $\Omega_n>\omega>\Omega_{n+1}$，
其中 $\Omega_j$ 是轨道上的电子公转角频率。
并证明大 $n$ 时 $\omega/\Omega_n\to1$。
\exsolution
玻尔能级写作 $E_n=-A/n^2$，$A>0$；
其第 $n$ 轨道公转角频率为
$\Omega_n=2A/(\hbar n^3)$。
相邻跃迁的光子角频率为
\[
 \omega=\frac{A}{\hbar}
 \left(\frac1{n^2}-\frac1{(n+1)^2}\right).
\]
对函数 $f(x)=-A/x^2$ 使用中值定理，
存在 $\xi\in(n,n+1)$ 使
$E_{n+1}-E_n=f'(\xi)=2A/\xi^3$，
因而
\[
 \boxed{\Omega_n>\omega>\Omega_{n+1}}.
\]
又 $\xi/n\to1$，得到
\[
 \boxed{\lim_{n\to\infty}
 \frac\omega{\Omega_n}=1}.
\]
\excheck 结论体现玻尔对应原理；
当 $n$ 大时相邻能级跃迁频率趋近经典公转频率。

\atomicproblem{21}{235}
质量为 $m$ 的粒子在中心势
$U(r)=kr^2/2$（$k>0$）中作圆周运动，
应用玻尔角动量量子化求允许的半径和能量。
\exsolution
圆周向心力条件
$mv^2/r=kr$，故公转角频率
$\Omega=v/r=\sqrt{k/m}$。
玻尔量子化
$L=mvr=mr^2\Omega=n\hbar$，
$n=1,2,\dots$，所以
\[
 \boxed{r_n=\sqrt{\frac{n\hbar}{m\Omega}}
 =\sqrt{\frac{n\hbar}{\sqrt{mk}}}}.
\]
机械能
$E_n=\tfrac12mv^2+\tfrac12kr_n^2
=kr_n^2$，于是
\[
 \boxed{E_n=n\hbar\Omega=
 n\hbar\sqrt{\frac km}}.
\]
\excheck 这是\emph{玻尔圆轨道}的模型，
不是一维量子谐振子精确本征值
$(n+\tfrac12)\hbar\Omega$。

\atomicproblem{22}{235}
对氢原子 H 和单电离氦离子 He$^+$，
求：(a) 第一玻尔轨道半径及电子速度；
(b) 基态动能与结合能；
(c) 电离电势、第一次激发电势、
共振谱线 $2\to1$ 波长。
\exsolution
固定核近似、核电荷数 $Z=1$ 或2，
\[
 r_{1,Z}=\frac{a_0}{Z},\quad
 v_{1,Z}=Z\alpha_{\rm fs}c,\quad
 K_{1,Z}=E_{b,Z}=13.606Z^2\,\mathrm{eV}.
\]
电离电势的数值（V）等于结合能的数值（eV）；
第一次激发所需能量为
$(1-1/4)E_b=3E_b/4$，
跃迁光子波长为 $\lambda_{21}=hc/(3E_b/4)$。
结果为
\[
\begin{array}{c|cc}
 &\mathrm H&\mathrm{He}^{+}\\\hline
 r_1\,(\mathrm{pm})&52.9&26.5\\
 v_1\,(10^6\mathrm{m/s})&2.18&4.36\\
 K_1=E_b\,(\mathrm{eV})&13.6&54.4\\
 U_{\mathrm{ion}}\,(\mathrm V)&13.6&54.4\\
 U_{\mathrm{exc},1}\,(\mathrm V)&10.2&40.8\\
 \lambda_{21}\,(\mathrm{nm})&121.6&30.4
\end{array}
\]
\excheck 原子 H 与类氢离子 He$^+$ 都只有一个电子，
同 $n$ 时 $r\propto Z^{-1}$、
$E\propto Z^2$、$\lambda\propto Z^{-2}$。

\atomicproblem{23}{235}
求单电离氦 He$^+$ 中电子沿第二玻尔轨道
运动的公转角频率。
\exsolution
类氢玻尔轨道半径
$r_n=a_0n^2/Z$，速度
$v_n=Z\alpha_{\rm fs}c/n$，
因此
\[
 \Omega_n=\frac{v_n}{r_n}
 =\frac{Z^2\alpha_{\rm fs}c}{a_0n^3}.
\]
对 He$^+$ 取 $Z=2,n=2$，
\[
 \boxed{\Omega_2\simeq2.07\times10^{16}\,\mathrm{s^{-1}}}.
\]
\excheck 求的是\emph{角频率}，
不是公转每秒圈数 $\Omega/(2\pi)$。

\atomicproblem{24}{235}
对类氢离子电子沿第 $n$ 个玻尔圆轨道运动，
求其轨道磁矩 $\boldsymbol\mu_n$、
轨道磁矩与轨道角动量之比；
给出第一轨道磁矩的模。
\exsolution
电子电荷为 $-e$，轨道等效电流
$I=-ev/(2\pi r)$，环路面积 $\pi r^2$，
故
\[
 \boldsymbol\mu=-\frac{evr}{2}\,\hat{\boldsymbol L}
 =-\frac e{2m_e}\boldsymbol L.
\]
按玻尔量子化 $|\boldsymbol L|=n\hbar$，
\[
 \boxed{\boldsymbol\mu_n=
 -n\frac{e\hbar}{2m_e}\,\hat{\boldsymbol L}},
 \qquad
 \boxed{\frac{|\boldsymbol\mu_n|}{|\boldsymbol L|}
 =\frac e{2m_e}},
\]
\[
 \boxed{|\boldsymbol\mu_1|=\mu_B=
 \frac{e\hbar}{2m_e}
 \simeq9.274\times10^{-24}\,\mathrm{J/T}}.
\]
\excheck 电子带负电，轨道磁矩矢量与
角动量反向；磁矩模是正值。

\atomicproblem{25}{235}
在玻尔模型中估算氢原子基态电子运动在
轨道中心产生的磁感应强度。
\exsolution
将电子轨道视为半径
$a_0=5.29\times10^{-11}\,\mathrm m$ 的圆电流环，
电流模为 $I=ev/(2\pi a_0)$，
环中心磁场为
\[
 B=\frac{\mu_0 I}{2a_0}
 =\frac{\mu_0 ev}{4\pi a_0^2}.
\]
用 $v=\alpha_{\rm fs}c\simeq2.19\times10^6\,\mathrm{m/s}$，
\[
 \boxed{B\simeq12.5\,\mathrm T
 =1.25\times10^5\,\mathrm G}.
\]
\excheck 这里估算的是经典玻尔\emph{轨道电流}磁场，
真实基态电子没有经典意义上的绕核轨道；
不应把该数值说成精确量子自旋/轨道局域磁场。

\atomicproblem{26}{235}
计算氢原子莱曼系、巴耳末系、帕邢系
在波长轴上覆盖的区间，并标出可见光区域。
\exsolution
里德伯公式
\[
 \frac1\lambda=R_H\left(\frac1{n_1^2}-\frac1{n_2^2}\right),
 \qquad n_2>n_1.
\]
对同一系，最长波长来自 $n_2=n_1+1$，
最短波长（系限）来自 $n_2\to\infty$，
故
\[
 \lambda_{\max}=
 \frac1{R_H[n_1^{-2}-(n_1+1)^{-2}]},
 \quad
 \lambda_{\min}=\frac{n_1^2}{R_H}.
\]
取近似 $R_H=1.097\times10^7\,\mathrm{m^{-1}}$：
\[
 \boxed{\begin{array}{c|cc}
 &\lambda_{\min}&\lambda_{\max}\\\hline
 \text{莱曼 }(n_1=1)&91.2\,\mathrm{nm}&121.6\,\mathrm{nm}\\
 \text{巴耳末 }(n_1=2)&364.6\,\mathrm{nm}&656.3\,\mathrm{nm}\\
 \text{帕邢 }(n_1=3)&820.4\,\mathrm{nm}&1875\,\mathrm{nm}
 \end{array}}
\]
\excheck 可见区约为 $380$--$760\,\mathrm{nm}$，
其中的氢谱线主要属于巴耳末系；
该表可直接据端点在波长轴绘制三个连续区间。

\atomicproblem{27}{235}
氢原子某条谱线的波数，
等于巴耳末系波长
$486.1\,\mathrm{nm}$ 与 $410.2\,\mathrm{nm}$
两条谱线波数之差。
求该未知谱线属于哪一系及其波长。
\exsolution
巴耳末 $486.1\,\mathrm{nm}$ 为 $4\to2$，
$410.2\,\mathrm{nm}$ 为 $6\to2$。
其波数之\emph{差}为
\[
 \left[\frac1{410.2\,\mathrm{nm}}-
 \frac1{486.1\,\mathrm{nm}}\right]
 =R_H\left(\frac1{4^2}-\frac1{6^2}\right).
\]
它对应从上能级6跃迁到下能级4，
即\emph{布拉开特系（Brackett）}的 $6\to4$ 谱线。
于是
\[
 \boxed{\lambda=\frac1{R_H(1/16-1/36)}
 \simeq2.63\,\mu\mathrm m}.
\]
\excheck 必须取两谱线\emph{波数}之差，
而不是直接拿两个波长相减。

\atomicproblem{28}{235}
对氢原子求：(a) 巴耳末系前三条谱线波长；
(b) 能分辨该系列前20条谱线的光谱仪，
至少应有多大的分辨本领 $\lambda/\delta\lambda$？
\exsolution
巴耳末系下能级为 $n_1=2$，
第 $k$ 条线对应上能级 $n_2=k+2$。
由里德伯公式，
\[
 \boxed{\lambda_1\simeq656\,\mathrm{nm},\quad
 \lambda_2\simeq486\,\mathrm{nm},\quad
 \lambda_3\simeq434\,\mathrm{nm}}.
\]
高序线间距不断缩小，最严格条件是区分
第20条（$22\to2$）与第21条（$23\to2$）。
故要求
\[
 \mathcal R\ge
 \frac{\lambda_{22}}{\lambda_{22}-\lambda_{23}},
 \quad
 \lambda_n=\frac1{R_H(1/4-1/n^2)}.
\]
直接计算
\[
 \boxed{\mathcal R_{\min}\simeq1.41\times10^3
 \approx1.5\times10^3}.
\]
\excheck 应比较最密集的一端；
不能只用第一条和第二条的大间隔估计。

\atomicproblem{29}{235}
宽度 $\ell=6.6\,\mathrm{mm}$ 的光栅
受到正入射的氢原子辐射。
巴耳末系的\emph{第50条谱线}在衍射角 $\theta$
附近恰达到瑞利分辨判据。
求 $\theta$，并对“第50条”序号作严格辨析。
\exsolution
令光栅实际参与干涉的缝数为 $N$，
衍射级次为 $m$。
光栅的瑞利分辨本领为
$\mathcal R=mN$，而正入射光栅
$m\lambda=d\sin\theta$、$N\simeq\ell/d$，
所以
\[
 \boxed{\sin\theta=
 \frac{\lambda^2}{\ell\,\delta\lambda}}.
\]
上式中的分子分母约去后实质为
$\sin\theta=\lambda^2/(\ell\delta\lambda)$。
严格按\emph{第50条}定义，
上能级 $n=52$，邻线是 $n=53$。
用里德伯公式
$\lambda_{52}\simeq365.047\,\mathrm{nm}$、
$\delta\lambda\simeq0.02021\,\mathrm{nm}$，
则
\[
 \boxed{\sin\theta\simeq0.999,\quad
 \theta\simeq87^\circ}.
\]
\textbf{原书答案约 $60^\circ$ 的来源：}
若把“第50条”\emph{直接当作上能级 $n=50$}，
则相邻谱线 $50\to2$、$51\to2$
给 $\sin\theta\simeq0.889$、$\theta\simeq63^\circ$，
即书末近似 $60^\circ$。
\excheck 第50条谱线与 $n=50$ 上能级
原则上差2；本题的原书题号解释需由独立学术复核决定，
不能悄然抹去两种编号造成的差异。

\atomicproblem{30}{236}
哪种元素的类氢谱线波长等于氢原子对应谱线的四分之一？
\exsolution
类氢离子（忽略约化质量细修正）的能级
$E_n\propto-Z^2/n^2$，
对应两能级间辐射的光子能量
$\Delta E\propto Z^2$。
因此对应波长
$\lambda_Z=\lambda_H/Z^2$。
由 $\lambda_Z=\lambda_H/4$ 得 $Z=2$，
所以
\[
 \boxed{\mathrm{He}^{+}\ \text{（单电离氦离子）}}.
\]
\excheck 普通氦原子有两个电子，
不是类氢单电子体系；题中对应的是 He$^+$。

\atomicproblem{31}{236}
氢原子被激发至主量子数 $n$ 的能级。
假定允许其在各个更低能级间辐射跃迁，
最多有多少条不同频率的谱线？
\exsolution
可参与级联跃迁的能级集合为
$1,2,\dots,n$。
每一对不同能级 $(i,j)$，$1\le i<j\le n$，
对应一次可能的向下跃迁，
故谱线种数为二元组合数
\[
 \boxed{N=\binom n2=\frac{n(n-1)}{2}}.
\]
\excheck 本题数的是\emph{整组被激发原子}
可能发出的谱线，不是单个原子一次级联必然发出的线数。

\atomicproblem{32}{236}
求氢原子\emph{吸收}谱在
$94.5$--$130.0\,\mathrm{nm}$ 区间内有哪些线。
初始原子均在基态。
\exsolution
基态吸收只能有 $1\to n$，$n\ge2$，
属于莱曼系，满足
\[
 \lambda_{1n}=
 \frac1{R_H(1-1/n^2)}.
\]
依次为
\[
 \boxed{\lambda_{12}\simeq121.6\,\mathrm{nm}},\quad
 \boxed{\lambda_{13}\simeq102.6\,\mathrm{nm}},\quad
 \boxed{\lambda_{14}\simeq97.3\,\mathrm{nm}},\quad
 \boxed{\lambda_{15}\simeq95.0\,\mathrm{nm}}.
\]
下一条 $1\to6$ 的波长约 $93.8\,\mathrm{nm}$，
已小于题给下界 $94.5\,\mathrm{nm}$。
\excheck 按底本题面区间应有\emph{四条}，
而底本书末只列前三条
$97.3,102.6,121.6\,\mathrm{nm}$，
漏掉仍在区间内的 $95.0\,\mathrm{nm}$；
已明确列为需校勘的答案缺项。
